\documentclass[11pt, a4paper]{article}

\usepackage[utf8]{inputenc}
\usepackage[T1]{fontenc}
\usepackage[french]{babel}
\usepackage{amsmath, amssymb}
\usepackage{geometry}
\usepackage{graphicx}
\usepackage{tabularx}
\usepackage{booktabs}
\usepackage{xcolor}
\usepackage{listings}
\usepackage{tikz}
\usepackage{hyperref}

\geometry{
    top=2cm,
    bottom=2cm,
    left=1.5cm,
    right=1.5cm
}

\lstset{
    basicstyle=\ttfamily\small,
    breaklines=true,
    frame=single,
    tabsize=4,
    showstringspaces=false,
    keywordstyle=\color{blue}\bfseries,
    commentstyle=\color{green!50!black},
    stringstyle=\color{red},
    numbers=left,
    numberstyle=\tiny\color{gray}
}

\newcommand{\imageplaceholder}[2]{
    \begin{center}
        \fbox{\begin{minipage}{#1}\centering\vspace{1cm}\textit{#2}\vspace{1cm}\end{minipage}}
    \end{center}
}

\newcommand{\corr}[1]{
    \vspace{0.3cm}
    \begin{center}
    \fcolorbox{blue}{blue!5}{
    \begin{minipage}{0.95\textwidth}
    \color{blue} \textbf{Correction \& Raisonnement :} \\
    #1
    \end{minipage}}
    \end{center}
    \vspace{0.3cm}
}

\begin{document}

% ==================== EN-TÊTE ====================
\noindent
\begin{tabularx}{\textwidth}{@{}X c r@{}}
\textbf{IN361} & \textbf{Session de juin 2025} & \textbf{Documents selon sujet} \\
\textbf{Java} & \textbf{Examen} & \textbf{Calculatrice selon sujet} \\
\end{tabularx}
\vspace{0.2cm}
\hrule
\vspace{0.4cm}
% ==================== EXERCICE 1 ====================
\begin{center}
    \textbf{\Large Exercice 1 (3 points)}
\end{center}
\vspace{-0.2cm}
\hrule
\vspace{0.3cm}
\textbf{Question :} Quel est le résultat de ce programme?

\corr{
\textbf{Résultat et Raisonnement :}
Le programme illustre les principes d'héritage et de polymorphisme (liaison dynamique).
\begin{itemize}
    \item \texttt{A a = new A(); a.affiche();} $\rightarrow$ \textbf{Je suis un A}
    \item \texttt{B b = new B(); b.affiche();} $\rightarrow$ \textbf{Je suis un A} (B hérite de A et ne redéfinit pas la méthode ).
    \item \texttt{C c = new C(); c.affiche();} $\rightarrow$ \textbf{Je suis un C} (C redéfinit \texttt{affiche()} ).
    \item \texttt{a = c; a.affiche();} $\rightarrow$ \textbf{Je suis un C} (Polymorphisme : \texttt{a} pointe vers une instance de \texttt{C} ).
    \item \texttt{D d = new D(); d.affiche();} $\rightarrow$ \textbf{Je suis un D} (D redéfinit \texttt{affiche()} ).
    \item \texttt{a = d; a.affiche();} $\rightarrow$ \textbf{Je suis un D} (\texttt{a} pointe vers \texttt{D} ).
    \item \texttt{c = d; c.affiche();} $\rightarrow$ \textbf{Je suis un D} (\texttt{c} pointe vers \texttt{D} ).
    \item \texttt{E e = new E(); e.affiche();} $\rightarrow$ \textbf{Je suis un A} (E hérite de B, qui hérite de A. Aucune ne redéfinit \texttt{affiche()} ).
    \item \texttt{a = e; a.affiche();} $\rightarrow$ \textbf{Je suis un A} (\texttt{a} pointe vers \texttt{E} ).
    \item \texttt{b = e; b.affiche();} $\rightarrow$ \textbf{Je suis un A} (\texttt{b} pointe vers \texttt{E} ).
    \item \texttt{F f = new F(); f.affiche();} $\rightarrow$ \textbf{Je suis un C} (F hérite de C qui possède sa propre méthode \texttt{affiche()} ).
    \item \texttt{a = f; a.affiche();} $\rightarrow$ \textbf{Je suis un C} (\texttt{a} pointe vers \texttt{F} ).
    \item \texttt{c = f; c.affiche();} $\rightarrow$ \textbf{Je suis un C} (\texttt{c} pointe vers \texttt{F} ).
\end{itemize}
}

% ==================== EXERCICE 2 ====================
\begin{center}
    \textbf{\Large Exercice 2 (4 points)}
\end{center}
\vspace{-0.2cm}
\hrule
\vspace{0.3cm}

\subsection*{Question 2.1 }
\corr{
\textbf{Raisonnement :}
Création d'une classe classique avec attributs privés, constructeur et accesseurs/mutateurs.
\begin{lstlisting}[language=Java]
public class TypPneu {
    private int diametre;
    private int largeur;
    private String marque;

    public TypPneu(int diametre, int largeur, String marque) {
        this.diametre = diametre;
        this.largeur = largeur;
        this.marque = marque;
    }
    
    // Getters et Setters
    public int getDiametre() { return diametre; }
    public void setDiametre(int diametre) { this.diametre = diametre; }
    
    public int getLargeur() { return largeur; }
    public void setLargeur(int largeur) { this.largeur = largeur; }
    
    public String getMarque() { return marque; }
    public void setMarque(String marque) { this.marque = marque; }
}
\end{lstlisting}
}

\subsection*{Question 2.2 }
\corr{
\textbf{Raisonnement :}
La méthode \texttt{equals(Object o)} doit être redéfinie pour comparer les valeurs et non les références mémoires.
\begin{lstlisting}[language=Java]
@Override
public boolean equals(Object obj) {
    if (this == obj) return true;
    if (obj == null || getClass() != obj.getClass()) return false;
    TypPneu typPneu = (TypPneu) obj;
    return diametre == typPneu.diametre && 
           largeur == typPneu.largeur && 
           marque.equals(typPneu.marque);
}
\end{lstlisting}
}

\subsection*{Question 2.3 }
\corr{
\textbf{Raisonnement :}
Création de la classe \texttt{TypVoiture} avec composition (elle contient un objet \texttt{TypPneu}).
\begin{lstlisting}[language=Java]
public class TypVoiture {
    private String marque;
    private String modele;
    private TypPneu pneu;

    public TypVoiture(String marque, String modele, TypPneu pneu) {
        this.marque = marque;
        this.modele = modele;
        this.pneu = pneu;
    }

    // Getters et Setters
    public String getMarque() { return marque; }
    public void setMarque(String marque) { this.marque = marque; }

    public String getModele() { return modele; }
    public void setModele(String modele) { this.modele = modele; }

    public TypPneu getTypPneu() { return pneu; }
    public void setTypPneu(TypPneu pneu) { this.pneu = pneu; }
}
\end{lstlisting}
}

\subsection*{Question 2.4 }
\corr{
\textbf{Raisonnement :}
Ajout d'une méthode validant les dimensions avant de remplacer le pneu.
\begin{lstlisting}[language=Java]
public void changementPneu(TypPneu tp) {
    if (this.pneu.getLargeur() == tp.getLargeur() && 
        this.pneu.getDiametre() == tp.getDiametre()) {
        this.pneu = tp;
    } else {
        throw new RuntimeException("Dimensions incompatibles");
    }
}
\end{lstlisting}
}

\subsection*{Question 2.5 }
\textbf{Erreur : } La méthode \texttt{changementPneu} est appelée sur la classe \texttt{TypVoiture} (\texttt{TypVoiture.changementPneu(p2);}) au lieu d'être appelée sur l'instance \texttt{v1}. Le code correct serait : \texttt{v1.changementPneu(p2);}

\subsection*{Question 2.6 }
\corr{
\textbf{Raisonnement :}
\texttt{==} compare les références mémoires. \texttt{equals()} compare le contenu.
\begin{itemize}
    \item \texttt{p1 == p2} $\rightarrow$ \textbf{false} (Différentes instances en mémoire).
    \item \texttt{p1 == p3} $\rightarrow$ \textbf{false} (Différentes instances).
    \item \texttt{p1.equals(p2)} $\rightarrow$ \textbf{false} (Marques différentes : "Michelin" $\neq$ "Dunlop" ).
    \item \texttt{p1.equals(p3)} $\rightarrow$ \textbf{false} (Dimensions différentes ).
    \item \texttt{p1 == p4} $\rightarrow$ \textbf{false} (\texttt{p4} pointe vers \texttt{p3} ).
    \item \texttt{p3 == p4} $\rightarrow$ \textbf{true} (Même référence ).
    \item \texttt{p3.equals(p4)} $\rightarrow$ \textbf{true} (Même référence, donc même contenu).
    \item \texttt{v1.getTypPneu() == v2.getTypPneu()} $\rightarrow$ \textbf{true} (Les deux voitures partagent l'objet \texttt{p1} ).
    \item \texttt{v1.getTypPneu() == v3.getTypPneu()} $\rightarrow$ \textbf{false} (\texttt{p1} $\neq$ \texttt{p3}).
    \item \texttt{v1.getTypPneu().equals(p2)} $\rightarrow$ \textbf{false} (Équivaut à \texttt{p1.equals(p2)}).
\end{itemize}
}

% ==================== EXERCICE 3 ====================
\begin{center}
    \textbf{\Large Exercice 3 (3 points)}
\end{center}
\vspace{-0.2cm}
\hrule
\vspace{0.3cm}

\subsection*{3.1 }
On crée un objet en utilisant le mot-clé \texttt{new} suivi de l'appel au constructeur de la classe (ex: \texttt{A a = new A();}).

\subsection*{3.2 }
En Java, le développeur ne supprime pas explicitement les objets. C'est le \textbf{Garbage Collector} (Ramasse-miettes) qui détruit automatiquement les objets en mémoire lorsqu'ils ne sont plus référencés.

\subsection*{3.3 }
Java est \textbf{à la fois compilé et interprété}. Le code source (.java) est d'abord compilé en un code intermédiaire appelé \textbf{Bytecode} (.class). Ce bytecode est ensuite interprété (et souvent compilé à la volée, JIT) par la Machine Virtuelle Java (JVM) lors de l'exécution.

\subsection*{3.4 }
Le mot-clé \texttt{final} pour une variable indique qu'elle est une \textbf{constante}. Une fois initialisée, sa valeur ne peut plus être modifiée.
\textit{Exemple :} \texttt{final double PI = 3.14159;} Si on tente de réassigner \texttt{PI = 3.14;}, le compilateur renverra une erreur.

\subsection*{3.5 }
\begin{itemize}
    \item \textbf{Classe abstraite :} Peut contenir des attributs (état), des méthodes concrètes et des constructeurs. Une classe ne peut hériter que d'une seule classe (abstraite ou non).
    \item \textbf{Interface :} Ne contient (traditionnellement) que des signatures de méthodes abstraites et des constantes (\texttt{public static final}). Une classe peut implémenter plusieurs interfaces.
\end{itemize}

% ==================== EXERCICE 4 ====================
\begin{center}
    \textbf{\Large Exercice 4 (3 points)}
\end{center}
\vspace{-0.2cm}
\hrule
\vspace{0.3cm}

\begin{lstlisting}[language=Java]
public abstract class Vehicule {
    protected String immatriculation; // Format AA-111-AA sur 9 caracteres
    protected int poids;
    protected int places;

    public Vehicule(String immatriculation, int poids, int places) {
        this.immatriculation = immatriculation;
        this.poids = poids;
        this.places = places;
    }
    // Getters et setters omis pour breveté
}

public class Voiture extends Vehicule {
    private int placesEnfants;

    public Voiture(String immatriculation, int poids, int places, int placesEnfants) {
        super(immatriculation, poids, places);
        this.placesEnfants = placesEnfants;
    }
}

public class Camion extends Vehicule {
    private int chargeUtile;

    public Camion(String immatriculation, int poids, int places, int chargeUtile) {
        super(immatriculation, poids, places);
        this.chargeUtile = chargeUtile;
    }
}

public class Moto extends Vehicule {
    private int cylindree;

    public Moto(String immatriculation, int poids, int places, int cylindree) {
        super(immatriculation, poids, places);
        this.cylindree = cylindree;
    }
}
\end{lstlisting}

% ==================== EXERCICE 5 ====================
\begin{center}
    \textbf{\Large Exercice 5 (4 points)}
\end{center}
\vspace{-0.2cm}
\hrule
\vspace{0.3cm}

\begin{lstlisting}[language=Java]
public class Categorie {
    private char nom; // ex: 'A', 'B'...
    private double prixJour;

    public Categorie(char nom, double prixJour) {
        this.nom = nom;
        this.prixJour = prixJour;
    }
    public double getPrixJour() { return prixJour; }
}

// Modifications dans Vehicule
public abstract class Vehicule {
    // ... autres attributs (immatriculation, poids, places)
    protected Categorie categorie;

    public Vehicule(String immat, int poids, int places, Categorie categorie) {
        this.immatriculation = immat;
        this.poids = poids;
        this.places = places;
        this.categorie = categorie;
    }

    public double getPrixLocation(int nbJour) {
        return this.categorie.getPrixJour() * nbJour;
    }
}

// Exemple de modification du constructeur (idem pour Moto et Camion)
public class Voiture extends Vehicule {
    private int placesEnfants;
    public Voiture(String immat, int poids, int places, Categorie cat, int placesEnfants) {
        super(immat, poids, places, cat);
        this.placesEnfants = placesEnfants;
    }
}

// Classe de test
public class TestLocation {
    public static void main(String[] args) {
        Categorie catA = new Categorie('A', 50.0);
        Voiture v1 = new Voiture("AB-123-CD", 1200, 5, catA, 2);
        System.out.println("Prix pour 3 jours: " + v1.getPrixLocation(3));
    }
}
\end{lstlisting}

% ==================== EXERCICE 6 ====================
\begin{center}
    \textbf{\Large Exercice 6 (3 points)}
\end{center}
\vspace{-0.2cm}
\hrule
\vspace{0.3cm}

\begin{lstlisting}[language=Java]
import java.util.ArrayList;
import java.util.List;

public class Parc {
    private List<Vehicule> vehicules;

    public Parc() {
        this.vehicules = new ArrayList<>();
    }

    public void addVehicule(Vehicule v) {
        for (Vehicule existant : vehicules) {
            if (existant.immatriculation.equals(v.immatriculation)) {
                throw new RuntimeException("Vehicule deja existant");
            }
        }
        vehicules.add(v);
    }

    public void removeVehicule(String immat) {
        Vehicule aSupprimer = null;
        for (Vehicule v : vehicules) {
            if (v.immatriculation.equals(immat)) {
                aSupprimer = v;
                break;
            }
        }
        if (aSupprimer == null) {
            throw new RuntimeException("Vehicule introuvable");
        }
        vehicules.remove(aSupprimer);
    }

    public void displayParc() {
        for (Vehicule v : vehicules) {
            System.out.println("Immat: " + v.immatriculation + 
                ", Prix jour: " + v.getPrixLocation(1));
            // L'ideal serait d'avoir une methode toString() dans Vehicule
        }
    }

    public double getPrixTotalLocation(int nbJour) {
        double total = 0;
        for (Vehicule v : vehicules) {
            total += v.getPrixLocation(nbJour);
        }
        return total;
    }
}
\end{lstlisting}

\end{document}
